% ------------------------------------------------------------
% Page 81
% ------------------------------------------------------------

La perspective d’un ravitaillement meilleur, avec les possibilités de culture et une plus grande
sécurité nous firent opter pour cette solution. Nous venions d’avoir notre 2\textsuperscript{o}
enfant lorsqu’en Août 1941 nous quittâmes la Côte d’Azur.

\vspace{0.45cm}

J’entrais en fonction peu après. Mon travail de Greffier m’occupait environ 10 (?) heures par
mois. Je consacrais le reste de mon temps à la propriété. Cela dura ainsi jusqu’au printemps
1943. C’est alors que j’appris que l’usine de Carburants Forestiers embauchait du personnel
et qu’elle allait bientôt démarrer.

\vspace{0.45cm}

Connaissant la comptabilité, je me présentais comme comptable. Je fus accepté par le Conseil
d’Administration qui acceptait que je prélève 2h par mois pour la tenue des audiences de la
justice de paix.

\vspace{0.45cm}

Je pris donc mes fonctions le 1\textsuperscript{er} Juillet 1943. L’usine n’était pas tout à fait achevée. On
constituait le stock de bois. Je fus, au cours de ce mois de juillet envoyé en stage à
Chapareillan en Isère, pendant une semaine.

\vspace{0.45cm}

Début Août 1943 l’usine était prête pour les essais. Le Directeur était un jeune ingénieur des
Arts et Métiers M. Georges Rolland, son adjoint un dessinateur projeteur M. Charles. Nous
avions à l’époque de 28 à 33 ans. J’étais le plus âgé. Nous sommes depuis devenus
d’excellents amis.

\vspace{0.45cm}

\underline{Le PDG adjoint M. Fouilloux m’avait dit que nous étions classés usine prioritaire, que nous aurions à livrer une partie de notre production aux Allemands et que cela dispensait tout le personnel du STO. Avantage fort appréciable.}

\vspace{0.55cm}

\underline{\textbf{L’attaque du maquis d’Aire de Côte en 1943}}

\vspace{0.35cm}

\underline{L’affaire} d’Aire de Côte s’était passée le 1\textsuperscript{er}. juillet, \textit{( il semble qu’il y ait erreur car ce
maquis fut attaqué en mai 43 avec 7 morts et 39 prisonniers déportés )} jour de mon
embauche. Plusieurs réfractaires eurent la chance de pouvoir s’échapper. L’un d’eux s’était
réfugié à Ayres. Il vint se faire embaucher après avoir pris contact avec le pasteur Robert. Je
l’inscrivis sous le nom de Guigues.

\vspace{0.45cm}

Pour toute embauche, il fallait remplir une fiche avec tous les renseignements d’identité
concernant le candidat et l’emploi qu’il devait occuper. Cette fiche était envoyée au STO à
Mende qui donnait son autorisation. Il nous était interdit d’embaucher sans autorisation du
STO.

\vspace{0.45cm}

Je me souviens du repas d’inauguration offert par la Direction. Il eut lieu un des premiers
dimanches d’Août 1943. Notre ami « Marius » qui assurait les liaisons et le ravitaillement des
bûcherons s’était arrangé pour nous faire faire un excellent repas dans l’entrée de l’usine. Il y
avait un garçon embauché comme chef de four sous le nom d’Alleyard ou Tallevard. J’avais
enregistré contre lui un PV de Gendarmerie pour une contravention bénigne ; il ne s’était pas
présenté à l’audience. Les gendarmes le recherchaient comme réfractaire au STO. Il en fût
averti et mon camarade Sylvain Rey, hôtelier à Meyrueis vint me dire : « Ce garçon est
